\documentclass[../main.tex]{subfiles}

\begin{document}

\section{Introduction}
\label{sec:intro}

\subsection{Purpose and Engagement Type}
\label{sec:engagement-purpose}

This report presents a targeted, evidence-based technical assessment of
safety, security, and privacy risks in the HDF5 software and data ecosystem.
Its purpose is to identify systemic risk patterns, explain the trust boundaries
at which those risks arise, and recommend mitigations for the organizations and
individuals that develop, distribute, deploy, or use HDF5.

The term \emph{audit} denotes a documented technical examination. This
engagement is not a compliance certification, attestation engagement,
penetration test, or exhaustive line-by-line source-code review. It expresses
no independent assurance opinion and does not certify that any HDF5 release,
application, file, or deployment is safe, secure, privacy-preserving,
conforming, or free from defects. Authorship and sponsorship are disclosed in
the front matter.

\subsection{Objectives and Intended Audience}
\label{sec:engagement-objectives}

The objectives of the assessment are to:

\begin{enumerate}
    \item identify important assets, interfaces, capabilities, trust boundaries,
          and failure modes across the HDF5 ecosystem;
    \item distinguish accidental safety and data-integrity failures from
          adversarial security failures and privacy harms arising during
          authorized operation;
    \item identify recurring root-cause classes rather than treating individual
          crashes, CVEs, disclosures, and interoperability failures as
          unrelated events;
    \item examine how risk propagates between the file format, core library,
          tools, extensions, bindings, applications, supply chain, and
          independent implementations;
    \item distinguish file-format conformance, acceptance by a particular
          implementation, and safe use in a particular deployment;
    \item recommend mitigations attributable to the actors able to implement
          and enforce them; and
    \item identify lessons applicable to comparable open-source scientific data
          ecosystems.
\end{enumerate}

The primary audience is The HDF Group's maintainers, security responders,
release engineers, and technical governance bodies. Additional audiences
include extension and plugin authors, binding and independent-implementation
maintainers, downstream application and distribution maintainers, operators of
data-ingestion services, research-infrastructure teams, data stewards, and
security and privacy reviewers.

Recommendations are assigned, where possible, to the layer able to enforce
them. A control failure at an application or deployment boundary is not
classified as a core-library defect merely because an HDF5 mechanism was
involved.

\subsection{Scope}
\label{sec:assessment-scope}

Component scope is not binary in this report. The assessment distinguishes the
following levels of coverage:

\begin{description}
    \item[Primary assessment objects.]
    These are the HDF5 on-disk format; the core C library, with emphasis on its
    handling of file-controlled metadata; HDF5 command-line tools; the capability,
    trust, and privacy characteristics of VOL connectors, filter plugins, and
    Virtual File Drivers as extension classes; and HDF5 build, release, and
    dependency configuration. The report states findings and mitigations for
    these objects.

    \item[Boundary case studies.]
    Selected downstream incidents and interoperability disputes are examined to
    determine how HDF5 capabilities interact with application authority,
    deployment policy, or independent implementations. The cases include model
    and dataset ingestion, external-resource resolution, plugin loading, and
    conformance disputes involving files produced by independent writers. The
    referenced projects are not audited as products.

    \item[Contextual references.]
    Other components are identified to describe ecosystem reach, dependency
    relationships, compatibility consequences, or potential exposure. A
    contextual reference does not constitute an assessment or assurance
    conclusion about that component.
\end{description}

Case studies and samples are selected purposively according to exposure to
untrusted input or ambient authority, evidence of a recurring failure class,
interoperability or archival consequence, and relevance to cross-cutting
mitigations. They are not a statistical sample and do not support prevalence or
frequency estimates.

In this report, \emph{ecosystem-wide} describes the breadth of surfaces
considered. It does not mean that every implementation, configuration, or
downstream product was examined exhaustively.

\paragraph{Out of scope.}
The following are outside the engagement:

\begin{itemize}
    \item HSDS source code, client behavior, REST endpoints, authentication and
          authorization, tenant isolation, deployment configuration, storage
          backends, and service availability;
    \item a complete safety, security, or privacy assessment of HDFView, Keras,
          TensorFlow, Hugging Face, PureHDF, or any other application, service,
          framework, or independent implementation cited as a case study or
          ecosystem dependency;
    \item general AI or machine-learning safety, model behavior, training,
          inference, and framework functionality unrelated to an HDF5
          integration boundary;
    \item detailed assessment of MPI semantics, collective-I/O correctness,
          concurrency, fault tolerance, performance, or parallel-filesystem
          security in Parallel HDF5;
    \item exhaustive review or testing of every HDF5 release, platform, build
          configuration, compiler, filter, connector, binding, application,
          package, or independent implementation;
    \item penetration testing of production services, exploit development, or
          proof that every defect or vulnerability class has been identified;
    \item assessment of a particular organization's production infrastructure,
          credentials, identity systems, network controls, or operational
          procedures, except as context for a cited incident or mitigation; and
    \item legal, regulatory, contractual, research-ethics, or functional-safety
          certification.
\end{itemize}

References to HDFView identify downstream exposure or reader behavior, not an
HDFView finding. References to machine-learning frameworks or services concern
only their handling of HDF5-controlled inputs. References to Parallel HDF5 are
limited to the configuration, dependency, extension, or trust-boundary
characteristics expressly discussed.

HSDS may appear in ecosystem diagrams for context, but no conclusion about
HSDS should be inferred.

\subsection{Assessment Dimensions}
\label{sec:assessment-dimensions}

\paragraph{Safety.}
Safety means operational robustness and protection of data integrity,
availability, recoverability, interoperability, and long-term accessibility in
the absence of a malicious actor. Relevant conditions include malformed,
truncated, inconsistent, or partially written files; resource exhaustion; I/O
failure; incompatible implementations; and foreseeable misuse. Safety in this
report does not mean regulated functional or physical safety.

\paragraph{Security.}
Security concerns the possibility that an adversary can compromise
confidentiality, integrity, or availability, or cause a process to exercise
authority contrary to the operator's intent. The principal threat scenario is
an HDF5 file, metadata field, external reference, plugin requirement, archive,
or model artifact controlled by an untrusted party and processed by an
application possessing filesystem, network, compute, plugin-loading, or
credential authority.

The assessment includes memory-safety defects, denial of service, resource
exhaustion, external-resource resolution, dynamic code loading, downstream
deserialization, and supply-chain compromise. A security consequence does not
require a defect in the HDF5 format or parser: an intentional HDF5 capability
can become an attack mechanism when a downstream application interprets
file-controlled instructions using authority unavailable to the file's
originator.

\paragraph{Privacy.}
Privacy concerns adverse consequences arising from the disclosure, retention,
inference, aggregation, correlation, or secondary use of information. Such
consequences may arise during authorized and otherwise correct operation,
without a security breach. The assessment includes scientific data,
descriptive and structural metadata, object names and attributes, provenance,
workflow artifacts, caches, temporary copies, external resources, and
information that becomes sensitive when datasets are combined.

The report evaluates privacy-exposure mechanisms and possible technical or
governance controls. It does not determine compliance with privacy statutes,
regulations, contracts, consent requirements, institutional policies, or
research-ethics obligations.

\subsection{Temporal Boundary and Evidence Basis}
\label{sec:evidence-boundary}

The evidence cutoff for this revision is August 27, 2026. Public evidence
available on or before that date is eligible for consideration. Later
developments are outside this revision unless expressly incorporated as
subsequent-event updates.

The historical analysis has no fixed lower date and includes earlier releases
where necessary to examine cited vulnerabilities, compatibility behavior, and
implementation changes. Because this is a historical, ecosystem-level
assessment, no single software baseline applies. A conclusion concerning one
release, commit, build configuration, platform, or third-party version is not
presumed to apply to others. Version-specific conclusions are bounded by the
artifact identifiers cited with them. Where an exact artifact is not identified,
a statement should be read as an architectural or historical class-level
observation, not as a release-specific conclusion.

The principal evidence consists of published HDF5 specifications and API
documentation; selected source code and build configurations; security policies,
advisories, and CVE records; issue and pull-request discussions; cited commits
and regression artifacts; and authoritative documentation or incident reports
from representative downstream and independent projects.

\subsection{Methodology}
\label{sec:assessment-methodology}

The engagement uses a qualitative, risk-informed technical-assessment method:

\begin{enumerate}
    \item \textbf{Surface and trust-boundary mapping.}
          The ecosystem is decomposed into format, parser, tool, extension,
          binding, application, supply-chain, independent-implementation, and
          data-lifecycle surfaces.

    \item \textbf{Asset and harm analysis.}
          The assessment identifies relevant data, metadata, code, credentials,
          compute and storage resources, availability, interoperability, and
          interpretability, together with plausible safety, security, and
          privacy consequences.

    \item \textbf{Document and artifact examination.}
          Specifications, documentation, public source artifacts, vulnerability
          records, fixes, issue records, build configuration, and representative
          ecosystem documentation are examined.

    \item \textbf{Root-cause and pattern analysis.}
          Defects and incidents are grouped by the earliest violated invariant
          or trust-boundary failure supported by the evidence, rather than only
          by crash type, CVE identifier, or affected application.

    \item \textbf{Conformance and interoperability analysis.}
          Acceptance by one HDF5 reader is not treated as proof that a file
          conforms to the applicable specification, and rejection by another
          reader is not treated as proof that it does not. Where implementations
          disagree, the analysis distinguishes the normative requirement,
          writer behavior, reader validation, and compatibility policy.

    \item \textbf{Mitigation analysis.}
          Controls are evaluated according to potential impact, exposure to the
          failure mode, breadth and recurrence, existing controls, evidence
          strength, urgency, implementation dependencies, and the actor able to
          enforce them. Controls addressing cross-cutting root causes or enabling
          prerequisite protections are distinguished from local remediations.
          Risk severity and implementation sequence are treated as separate
          judgments.
\end{enumerate}

For incidents involving multiple components, the analysis distinguishes the
HDF5 mechanism, the authority supplied by the host application, the failed
control or policy boundary, and the resulting impact. The presence of HDF5 in
an exploit chain is not, without additional evidence, classified as a parser or
file-format defect.

The report distinguishes, where practicable, among:

\begin{description}
    \item[Documented observation:] directly supported by a cited specification,
          source artifact, test, advisory, incident record, or maintainer
          diagnosis; the applicable discussion distinguishes source-reported
          evidence from independent reproduction;
    \item[Reasoned inference:] a conclusion drawn from observations, with
          material assumptions and uncertainty identified; and
    \item[Recommendation:] a proposed control to be validated through
          implementation and acceptance evidence.
\end{description}

Unless a finding expressly records an independent reproduction or test, it
should be understood as the result of document, artifact, or architecture
examination rather than independent dynamic verification.

\subsection{Limitations and Assurance}
\label{sec:assessment-limitations}

This assessment is time-bounded and based predominantly on public evidence. It
is not an exhaustive review of the HDF5 source tree, every file-format feature,
every configuration, or every ecosystem component. Representative sampling
supports conclusions about observed mechanisms and recurring classes; it does
not establish their frequency throughout the ecosystem.

Public CVE and issue records may be incomplete, duplicated, disputed, or
subsequently corrected. Deployment-specific likelihood and impact depend on
application privileges, isolation, resource limits, enabled features, data
sensitivity, and provenance. Proposed mitigations are evaluated analytically
unless an implementation and verification procedure is expressly cited.

This report provides no formal assurance or certification opinion. Its
conclusions are bounded by the assessment objects, methods, evidence, and
cutoff described above. The absence of a reported finding does not establish
the absence of defects, vulnerabilities, safety failures, or privacy risk.

Mitigation priorities express the authors' risk-informed judgment. They become
auditable commitments only when paired with an owner, affected scope, target
horizon, implementation status, acceptance criteria, and residual-risk
decision.

\subsection{Privacy Context}
\label{sec:privacy-context}

The National Institute of Standards and Technology (NIST) defines privacy as “a condition that safeguards 
human autonomy and dignity through various means, including confidentiality, predictability, manageability, 
and disassociability.” Historically, privacy was interpreted as preventing unauthorized access to or 
disclosure of Personally Identifiable Information (PII). However, NIST emphasizes that contemporary 
privacy risks extend far beyond traditional notions of confidentiality. Modern privacy concerns increasingly 
arise from inference, aggregation, behavioral profiling, secondary use of data, context changes, and 
surveillance effects~\cite{nist_privacy_framework_2020}. NIST emphasizes that privacy risks are distinct from cybersecurity risks and 
may emerge even when systems are secure, lawful, and functioning exactly as intended~\cite{brooks2017privacy}. 
HDF5 inherits many of the same classes of privacy risks identified by NIST, although these risks manifest 
differently in scientific and engineering environments. 

HDF5 was designed to maximize data shareability, extensibility, portability, and long-term usability, 
and doesn’t enforce data privacy or provide data protection mechanism. Its rich metadata model, 
self-describing structure, plugin architecture, and interoperability layers can expose information 
extending beyond the raw data stored within the file. Sensitive information in HDF5 ecosystems may emerge 
through HDF5 internal and user-defined metadata, for example, HDF5 attributes, object names, 
file organization, and objects properties. It may be exposed during HDF5 data lifecycle through 
provenance records, workflow artifacts, and correlations across multiple datasets and files. 

A common misconception in scientific computing is that privacy concerns are resolved once 
data access is restricted or datasets are anonymized. However, privacy risks often persist even 
when traditional security mechanisms are correctly implemented. Restricting access does not prevent 
inference from metadata. Anonymization may remove direct identifiers while leaving contextual 
information sufficient for re-identification. Similarly, data reduction or aggregation does not 
eliminate the possibility that sensitive information may be reconstructed through correlations 
across files, workflows, or external repositories. In modern scientific ecosystems, privacy 
exposure frequently results from accumulation of data and metadata over time rather than disclosure 
of a single sensitive piece of raw data or metadata.

The privacy risks are amplified by current scientific practice of data openness and shareability.  
HDF5 datasets increasingly move across organization, systems, cloud platforms, workflows, and public 
and private repositories. Data are routinely repurposed beyond their original intent and analyzed by 
users who were not involved in their creation. Machine learning systems further amplify these concerns 
by discovering hidden relationships across data stored in HDF5 and enabling new forms of inference 
that may not have been anticipated when data were originally collected. As a result, data stored in 
the HDF5 files may experience evolving privacy characteristics through its lifecycle.

Therefore, privacy exposure in HDF5 is not limited to unauthorized access, data breaches,
unsafe use, or malicious behavior; it is distinct from, but can overlap with, data security and safety.
By design HDF5 file format and library do not guarantee data privacy. In addition, privacy risks may 
emerge during legitimate scientific workflows, through data sharing, reuse, and analysis.
In this document we audit privacy assets and different privacy exposure mechanisms (or surfaces) 
in HDF5 with the goal to raise awareness about privacy risks when using HDF5. 



\section{Assessment Surfaces}
\label{sec:assessment-surfaces}

Figure~\ref{fig:hdf5-assessment-surfaces} maps the principal components and
interfaces considered when identifying HDF5 risk. Inclusion in the figure
provides ecosystem context and does not override the scope levels and
limitations defined in Section~\ref{sec:assessment-scope}. In particular, HSDS
is shown for context but remains outside the engagement.

\begin{figure}[h!]
    \centering
    \includegraphics[width=\textwidth]{images/HDF5_Surfaces.png}
    \caption{HDF5 assessment surfaces (interfaces: rounded rectangles)}
    \label{fig:hdf5-assessment-surfaces}
\end{figure}

\end{document}
